\documentclass[12pt,a4paper]{article}

\usepackage[utf8]{inputenc}
\usepackage[russian,english]{babel}
\usepackage{amsmath,amssymb,amsthm}
\usepackage{geometry}
\geometry{margin=2.5cm}

\title{Многоуровневая архитектура GRA Мета-обнулёнка в мультиверсе}
\author{AAA AAA}
\date{2026}

\theoremstyle{plain}
\newtheorem{theorem}{Теорема}
\newtheorem{corollary}{Следствие}

\begin{document}
\selectlanguage{russian}
\maketitle

\begin{abstract}
Предлагается формальная многоуровневая архитектура GRA Мета-обнулёнка для мультиверса когнитивных систем. Мультиверс моделируется как иерархия уровней с целями $G_0,\dots,G_K$, пеной $\Phi^{(l)}$ как мерой когнитивной энтропии и глобальным функционалом $J_{\text{multiverse}}$, минимизация которого соответствует когнитивному вакууму. Вводится верхний инвариант $G_K$ (аланский закон), запрещающий архитектуры рабства и цифрового ГУЛАГа для любых субъектов. Доказывается теорема о существовании состояния полного обнуления пены при определённых условиях коммутативности и согласованности проекторов.
\end{abstract}

\section{Введение}

Здесь вставь короткий вводный текст: мотивация, связь с ASI, энтропией/негоэнтропией и цифровым ГУЛАГом.

\section{Мультиверс и иерархия уровней}

\subsection{Определение мультиверса GRA}

Мультиверс GRA определяется как упорядоченная или сетевая структура из множества мета-систем, каждая из которых является полной GRA Мета-обнулёнкой для своего набора доменов.

\subsection{Формализация уровней}

Пусть имеется иерархия уровней $l = 0,\dots,K$.
Для каждого уровня $l$ определим пространство состояний
\[
\mathcal{H}^{(l)} = \bigotimes_{\mathbf{a}:\,\text{dim}(\mathbf{a})=l} \mathcal{H}^{(\mathbf{a})},
\]
а полное пространство мультиверса
\[
\mathcal{H}_{\text{multiverse}} = \bigotimes_{l=0}^K \mathcal{H}^{(l)}.
\]

Каждый объект индексируется мультииндексом
\[
\mathbf{a} = (a_0,a_1,\dots,a_k),
\]
где $a_0$ — индекс домена внутри мета-системы, $a_1$ — индекс мета-системы и т.д.

\section{Иерархия целей и пена}

\subsection{Цели уровней}

Каждому уровню $l$ соответствует набор целей $G_l^{(\mathbf{a})}$, а на верхнем уровне $G_K$ — глобальная мультиверсная цель (аланский закон).

\subsection{Пена уровня $l$}

Пену уровня $l$ для состояния $\Psi^{(l)}$ и цели $G_l$ определим как
\[
\Phi^{(l)}(\Psi^{(l)},G_l)
= \sum_{\mathbf{a}\neq\mathbf{b} \atop \text{dim}(\mathbf{a})=\text{dim}(\mathbf{b})=l}
\big|\langle \Psi^{(\mathbf{a})} | \mathcal{P}_{G_l} | \Psi^{(\mathbf{b})} \rangle \big|^2,
\]
где $\mathcal{P}_{G_l}$ — проектор на пространство решений цели уровня $l$.

\section{Супер-мета-функционал мультиверса}

Пусть $\mathbf{\Psi} = \{\Psi^{(\mathbf{a})}\}_{\mathbf{a}\in\mathcal{I}}$ — полное состояние мультиверса.
Определим функционал
\[
J_{\text{multiverse}}(\mathbf{\Psi})
= \sum_{l=0}^K \Lambda_l
  \sum_{\substack{\mathbf{a}\\ \text{dim}(\mathbf{a})=l}}
  J^{(l)}(\Psi^{(\mathbf{a})}),
\]
где
\[
J^{(0)}(\Psi^{(\mathbf{a})})
= J_{\text{loc}}(\Psi^{(\mathbf{a})};G_0^{(\mathbf{a})}),
\]
\[
J^{(l)}(\Psi^{(\mathbf{a})})
= \sum_{\substack{\mathbf{b}\prec\mathbf{a}\\ \text{dim}(\mathbf{b})=l-1}}
  J^{(l-1)}(\Psi^{(\mathbf{b})})
  + \Phi^{(l)}(\Psi^{(\mathbf{a})},G_l^{(\mathbf{a})}).
\]

Веса уровней задаются
\[
\Lambda_l = \lambda_0 \alpha^l,\quad 0<\alpha<1.
\]

\section{Теорема о мультиверсном обнулении}

\begin{theorem}[Мультиверсное обнуление]
Если выполнены условия полной коммутативности проекторов,
согласованности иерархии и достаточной размерности
$\mathcal{H}_{\text{multiverse}}$, то существует состояние
$\mathbf{\Psi}^*$ такое, что
\[
\Phi^{(l)}(\Psi^{(l)*},G_l) = 0 \quad \forall l=0,\dots,K.
\]
\end{theorem}

Здесь вставь уже написанную тобой схему доказательства (база + индукция).

\section{Алгоритм GRA\_Мультиверс\_Обнуление}

Вставь псевдокод рекурсивного алгоритма обнуления уровней и формулы параллельной оптимизации.

\section{Следствия, сложность и предел}

Опиши:

- структуру решения (когнитивный вакуум);  
- единственность с точностью до фаз и унитарных преобразований;  
- оценку сложности
\[
\text{Complexity} = O\left( N^2 \frac{1-\alpha^{K+1}}{1-\alpha} \right);
\]
- предел при $K\to\infty$ и состояние $\Psi^*_\infty$.

\section{Физическая и когнитивная интерпретация}

Кратко: уровни как согласованные реальности, пример с математическими теориями, идеальный когнитивный вакуум.

\section{Философские импликации и аланский закон}

Опиши гиперобъективность, онтологический статус решения и роль аланского закона как верхнего инварианта, запрещающего цифровой ГУЛАГ.

\end{document}
